\chapter{Derivation of the Finite-Element Stiffness Method for Pin-Jointed Trusses}
\label{app:fem}

This appendix derives from first principles the finite-element
formulation used throughout the thesis for pin-jointed steel
trusses. The target audience is a civil-engineering reader who
has seen structural analysis but wants to see every step that
maps the continuum statement of equilibrium onto the matrix
expression $\mathbf{K}\mathbf{u}=\mathbf{F}$ that our Python
solver (\texttt{src/fem/solver.py}) actually evaluates. The
derivation is kept dimension-agnostic so the same formulae cover
the \num{10}-bar planar, \num{25}-bar spatial, \num{72}-bar spatial
and \num{200}-bar planar benchmarks of \cref{chap:results}.

\section{Strong Form: Equilibrium of a Pin-Jointed Truss}
\label{app:fem:strong}

A pin-jointed truss is an assembly of straight two-force members.
Every joint is an ideal frictionless pin; moments vanish at each
joint; and every member therefore carries only an axial force
(pure tension or pure compression). The three physical hypotheses
that underpin the entire chapter are:
\begin{description}
  \item[H1 — pure axiality.] Every member force lies along its axis.
  \item[H2 — small strain.] Deformations are infinitesimal; geometry
    is unchanged when writing equilibrium (first-order theory).
  \item[H3 — linear elasticity.] Each member obeys Hooke's law with
    a constant Young's modulus~$E$.
\end{description}

Under these hypotheses, static equilibrium of the joints is a
purely linear-algebraic statement. Let the truss have $N$ joints.
For any joint $p$ the vector sum of external load $\mathbf{f}_p$
and the internal member forces $\mathbf{t}_p^{(e)}$ (contributions
from every member $e$ incident at $p$) vanishes:
\begin{equation}
  \sum_{e\,\ni\,p}\mathbf{t}_p^{(e)}\;+\;\mathbf{f}_p
  \;=\;\mathbf{0},
  \qquad p=1,\dotsc,N.
  \label{eq:strong_equilibrium}
\end{equation}
Stacking \cref{eq:strong_equilibrium} across every joint and every
coordinate direction produces a system of $Nd$ scalar equations,
where $d=2$ for a planar truss and $d=3$ for a spatial truss. The
goal of the finite-element method is to re-express the left-hand
side as $\mathbf{K}\mathbf{u}$, where $\mathbf{K}$ is a structural
stiffness matrix depending only on geometry and material and
$\mathbf{u}$ is the unknown nodal-displacement vector. Doing so
reduces \cref{eq:strong_equilibrium} to
$\mathbf{K}\mathbf{u}=\mathbf{F}$, which is what the code solves.

\section{Bar Kinematics: Axial Strain from Nodal Displacements}
\label{app:fem:kinematics}

Consider a single bar $e$ connecting joints $i$ and $j$, with
reference coordinates
$\mathbf{x}_i,\mathbf{x}_j\in\mathbb{R}^d$. Its undeformed length is
\begin{equation}
  L^{(e)} = \|\mathbf{x}_j-\mathbf{x}_i\|_2
\end{equation}
and its undeformed axial direction is the unit vector
\begin{equation}
  \hat{\mathbf{n}}^{(e)}
  = \frac{\mathbf{x}_j-\mathbf{x}_i}{L^{(e)}}.
  \label{eq:direction_cosines}
\end{equation}
In 2D, $\hat{\mathbf{n}}^{(e)}=(\cos\theta,\sin\theta)^\top$; in 3D,
$\hat{\mathbf{n}}^{(e)}=(l,m,n)^\top$ is the classical triple of
\emph{direction cosines}. The implementation caches both $L^{(e)}$
and $\hat{\mathbf{n}}^{(e)}$ at element construction time; see
\texttt{TrussElement.length} and \texttt{TrussElement.direction} in
\texttt{src/fem/truss\_element.py}.

Let $\mathbf{u}_i,\mathbf{u}_j\in\mathbb{R}^d$ be the unknown
displacement of the two end joints. Under \textbf{H2} the deformed
length is
\begin{equation}
  L_{\mathrm{def}}
  = \bigl\|(\mathbf{x}_j+\mathbf{u}_j)-(\mathbf{x}_i+\mathbf{u}_i)\bigr\|
  = L^{(e)}\bigl(1+\varepsilon^{(e)}\bigr)+\mathcal{O}(\|\mathbf{u}\|^2),
\end{equation}
where the axial engineering strain is the component of the relative
displacement along the bar:
\begin{equation}
  \varepsilon^{(e)}
  \;=\;\frac{\Delta L^{(e)}}{L^{(e)}}
  \;=\;\frac{\hat{\mathbf{n}}^{(e)}\cdot(\mathbf{u}_j-\mathbf{u}_i)}
             {L^{(e)}}.
  \label{eq:axial_strain}
\end{equation}
\Cref{eq:axial_strain} is exact to first order in $\|\mathbf{u}\|$
and consistent with H2.

It is convenient to pack the four (or six, in 3D) end-displacement
components into a single \emph{element nodal vector}
$\mathbf{u}^{(e)}\in\mathbb{R}^{2d}$:
\begin{equation}
  \mathbf{u}^{(e)}
  \;=\;
  \begin{pmatrix}\mathbf{u}_i\\\mathbf{u}_j\end{pmatrix}
  \in\mathbb{R}^{2d}.
\end{equation}
Then \cref{eq:axial_strain} can be written as a single
row-vector-times-column-vector contraction:
\begin{equation}
  \varepsilon^{(e)}
  \;=\;\mathbf{B}^{(e)}\mathbf{u}^{(e)},
  \qquad
  \mathbf{B}^{(e)}
  \;=\;\frac{1}{L^{(e)}}
  \begin{pmatrix}-\hat{\mathbf{n}}^{(e)\top}
                   &\hat{\mathbf{n}}^{(e)\top}\end{pmatrix}
  \in\mathbb{R}^{1\times 2d}.
  \label{eq:strain_displacement}
\end{equation}
$\mathbf{B}^{(e)}$ is the \emph{strain-displacement matrix} of the
element. The sign pattern $(-\hat{\mathbf{n}},+\hat{\mathbf{n}})$
reflects that pulling $j$ away from $i$ elongates the bar.

\section{Linear-Elastic Constitutive Law and Internal Force}
\label{app:fem:constitutive}

By H3, Hooke's law in its one-dimensional form relates axial
stress $\sigma^{(e)}$ to axial strain through the material
Young's modulus $E$:
\begin{equation}
  \sigma^{(e)} \;=\; E\,\varepsilon^{(e)}.
  \label{eq:hooke1d}
\end{equation}
The internal axial \emph{force} in the bar is the stress times
the cross-sectional area $A^{(e)}$:
\begin{equation}
  N^{(e)} \;=\; A^{(e)}\sigma^{(e)}
           \;=\; EA^{(e)}\varepsilon^{(e)}
           \;=\; \frac{EA^{(e)}}{L^{(e)}}\,\Delta L^{(e)}.
  \label{eq:axial_force}
\end{equation}
$N^{(e)}>0$ is tension, $N^{(e)}<0$ is compression. The quantity
$k_{\mathrm{ax}}^{(e)}\equiv EA^{(e)}/L^{(e)}$ is the bar's
\emph{axial stiffness} and plays the role of the spring constant
for this one-dimensional sub-problem.

\section{Element Stiffness Matrix via Virtual Work}
\label{app:fem:elem_stiffness}

The element-level stiffness matrix $\mathbf{k}^{(e)}$ is the
$2d\times 2d$ symmetric matrix that maps the bar's end
displacements $\mathbf{u}^{(e)}$ to the end forces
$\mathbf{r}^{(e)}$ the bar exerts on the two end nodes:
$\mathbf{r}^{(e)}=\mathbf{k}^{(e)}\mathbf{u}^{(e)}$. We derive
it via the \emph{principle of virtual work}, which for a single
elastic bar of length $L^{(e)}$ and area $A^{(e)}$ reads
\begin{equation}
  \int_0^{L^{(e)}}\!\!\sigma^{(e)}\,\delta\varepsilon^{(e)}\,A^{(e)}\,\mathrm{d}s
  \;=\;\mathbf{r}^{(e)\top}\delta\mathbf{u}^{(e)}
  \qquad\forall\,\delta\mathbf{u}^{(e)}.
  \label{eq:pvw}
\end{equation}
Both the stress and the strain are constant along the bar under
H1--H3, so the left-hand integral collapses:
\begin{align}
  \int_0^{L^{(e)}}\sigma^{(e)}\delta\varepsilon^{(e)}A^{(e)}\mathrm{d}s
  &= A^{(e)}L^{(e)}\,\sigma^{(e)}\delta\varepsilon^{(e)}
  \nonumber\\
  &= A^{(e)}L^{(e)}\,E\,\varepsilon^{(e)}\delta\varepsilon^{(e)}
  \nonumber\\
  &= A^{(e)}L^{(e)}\,E\,
     \bigl(\mathbf{B}^{(e)}\mathbf{u}^{(e)}\bigr)\!
     \bigl(\mathbf{B}^{(e)}\delta\mathbf{u}^{(e)}\bigr)
  \nonumber\\
  &= \delta\mathbf{u}^{(e)\top}\!
     \bigl(EA^{(e)}L^{(e)}\mathbf{B}^{(e)\top}\mathbf{B}^{(e)}\bigr)
     \mathbf{u}^{(e)}.
\end{align}
Substituting \cref{eq:strain_displacement} for
$\mathbf{B}^{(e)}$ (which contains the factor $1/L^{(e)}$ twice,
once from $\mathbf{B}^{(e)}$ and once from $\mathbf{B}^{(e)\top}$):
\begin{equation}
  EA^{(e)}L^{(e)}\mathbf{B}^{(e)\top}\mathbf{B}^{(e)}
  \;=\;\frac{EA^{(e)}}{L^{(e)}}
  \begin{pmatrix}-\hat{\mathbf{n}}\\+\hat{\mathbf{n}}\end{pmatrix}
  \begin{pmatrix}-\hat{\mathbf{n}}^\top & +\hat{\mathbf{n}}^\top\end{pmatrix}.
\end{equation}
Because \cref{eq:pvw} must hold for \emph{every} admissible
virtual displacement $\delta\mathbf{u}^{(e)}$, the coefficient
matrix is $\mathbf{k}^{(e)}$:
\begin{equation}
  \boxed{\;
  \mathbf{k}^{(e)}
  \;=\;\frac{EA^{(e)}}{L^{(e)}}
  \begin{pmatrix}
    +\hat{\mathbf{n}}\hat{\mathbf{n}}^\top & -\hat{\mathbf{n}}\hat{\mathbf{n}}^\top\\
    -\hat{\mathbf{n}}\hat{\mathbf{n}}^\top & +\hat{\mathbf{n}}\hat{\mathbf{n}}^\top
  \end{pmatrix}
  \;=\;\frac{EA^{(e)}}{L^{(e)}}\,
       \mathbf{B}^{(e)\top}\mathbf{B}^{(e)}\cdot L^{(e)\,2}\,
  }\;
  \label{eq:element_stiffness}
\end{equation}
where the last equality follows from recognising that writing
$\mathbf{B}^{(e)}$ without the $1/L^{(e)}$ factor turns it into
the rank-1 row vector
$\tilde{\mathbf{B}}^{(e)}\equiv L^{(e)}\mathbf{B}^{(e)}=
(-\hat{\mathbf{n}}^\top,\,\hat{\mathbf{n}}^\top)$ and the
coefficient is $(EA^{(e)}/L^{(e)})\,\tilde{\mathbf{B}}^{(e)\top}
\tilde{\mathbf{B}}^{(e)}$. In the code, the compact form
$\mathbf{k}^{(e)}=(EA/L)\,\mathbf{B}^\top\mathbf{B}$ with
$\mathbf{B}=(-\hat{\mathbf{n}}^\top,\,\hat{\mathbf{n}}^\top)$ is
what appears on \texttt{truss\_element.py:152}.

\paragraph{Properties of $\mathbf{k}^{(e)}$.} The element
stiffness matrix has three properties that every implementation
test should verify:
\begin{enumerate}
  \item \textbf{Symmetry.} $\mathbf{k}^{(e)}=\mathbf{k}^{(e)\top}$,
    because $\mathbf{B}^{(e)\top}\mathbf{B}^{(e)}$ is symmetric.
    This is physical: by Maxwell--Betti reciprocity, the force at
    joint $i$ due to a unit displacement at $j$ equals the force
    at $j$ due to a unit displacement at $i$.
  \item \textbf{Positive semidefiniteness.}
    $\mathbf{k}^{(e)}\succeq 0$, with one non-zero eigenvalue
    $2\,EA^{(e)}/L^{(e)}$ and $2d-1$ zero eigenvalues. The non-zero
    mode is axial stretch; the $2d-1$ zero modes are the rigid-body
    motions (translation in each of $d$ directions and
    $d(d-1)/2$ rigid rotations) and the purely transverse
    relative displacements which a pin-jointed bar does not resist.
  \item \textbf{Rank.} $\operatorname{rank}(\mathbf{k}^{(e)})=1$.
    Consequently, a truss assembled from $M$ bars has a global
    stiffness matrix whose rank cannot exceed $M$ before boundary
    conditions are applied — the assembled $\mathbf{K}$ will
    therefore in general be singular and demand support fixing
    (\cref{app:fem:bcs}).
\end{enumerate}

\section{Assembly: Summation over Elements}
\label{app:fem:assembly}

Assembly is the step that maps element contributions into the
global stiffness matrix. Number the joints $1,\dotsc,N$ and
associate with joint $p$ the $d$ consecutive global degrees of
freedom
$\{(p-1)d+1,(p-1)d+2,\dotsc,pd\}$. For an element connecting
joints $i,j$, the local-to-global DOF index map is the
$2d$-tuple
\begin{equation}
  \mathcal{I}^{(e)}
  =\bigl((i-1)d+1,\dotsc,id,\;(j-1)d+1,\dotsc,jd\bigr).
\end{equation}
The global stiffness matrix $\mathbf{K}\in\mathbb{R}^{Nd\times Nd}$
is assembled by scattering each element contribution onto the
rows and columns indexed by $\mathcal{I}^{(e)}$:
\begin{equation}
  \mathbf{K}[\mathcal{I}^{(e)},\mathcal{I}^{(e)}]
  \;\mathrel{+}=\;\mathbf{k}^{(e)},
  \qquad e=1,\dotsc,M.
  \label{eq:global_assembly}
\end{equation}
In our code (\texttt{src/fem/assembly.py:114}) the scatter is
implemented with \texttt{np.ix\_(dofs,dofs)} so the inner loop is
one line. $\mathbf{K}$ inherits three properties directly from
the element contributions:
\begin{itemize}
  \item symmetric ($\mathbf{k}^{(e)}=\mathbf{k}^{(e)\top}\;\forall\,e$);
  \item sparse (each bar scatters into only $(2d)^2$ entries,
        so average fill is $\mathcal{O}(M/N^2)$ which is tiny for
        triangulated trusses);
  \item positive semidefinite, with null space containing the
        rigid-body motions of the whole structure
        ($d$ translations and $d(d-1)/2$ rotations, giving a
        null-space dimension of $d(d+1)/2$).
\end{itemize}

\paragraph{Load vector.} The external loading is applied at the
joints. Defining $\mathbf{f}_p\in\mathbb{R}^d$ as the external
force at joint $p$, the global load vector is simply
\begin{equation}
  \mathbf{F}\;=\;
  \bigl(\mathbf{f}_1^\top,\mathbf{f}_2^\top,\dotsc,\mathbf{f}_N^\top\bigr)^\top
  \;\in\;\mathbb{R}^{Nd}.
\end{equation}
Self-weight is included by adding half the weight of each
incident member to each end joint, a lumped-mass distribution
that is exact under H1--H3 because bar forces are assumed
uniform.

\section{Boundary Conditions: Partition Method}
\label{app:fem:bcs}

Before boundary conditions the unassembled system
\begin{equation}
  \mathbf{K}\mathbf{u}\;=\;\mathbf{F}
\end{equation}
has no unique solution: $\mathbf{K}$ has a null space of dimension
$d(d+1)/2$ (see \cref{app:fem:assembly}), so for any particular
solution $\mathbf{u}^\ast$ we may add any rigid-body motion and
still satisfy the equation. Support conditions restore
uniqueness. Let
\begin{equation}
  \mathcal{F}=\{\text{free DOFs}\},\qquad
  \mathcal{S}=\{\text{prescribed DOFs}\},\qquad
  \mathcal{F}\cup\mathcal{S}=\{1,\dotsc,Nd\},\quad
  \mathcal{F}\cap\mathcal{S}=\varnothing.
\end{equation}
Partitioning $\mathbf{u}$ and $\mathbf{F}$ accordingly and
writing $\mathbf{K}$ in block form
\begin{equation}
  \begin{pmatrix}\mathbf{K}_{\mathcal{F}\mathcal{F}}&\mathbf{K}_{\mathcal{F}\mathcal{S}}\\
                 \mathbf{K}_{\mathcal{S}\mathcal{F}}&\mathbf{K}_{\mathcal{S}\mathcal{S}}\end{pmatrix}
  \begin{pmatrix}\mathbf{u}_{\mathcal{F}}\\\mathbf{u}_{\mathcal{S}}\end{pmatrix}
  \;=\;
  \begin{pmatrix}\mathbf{F}_{\mathcal{F}}\\\mathbf{F}_{\mathcal{S}}\end{pmatrix},
  \label{eq:partition_system}
\end{equation}
and using the prescribed support displacements
$\mathbf{u}_{\mathcal{S}}$ (commonly zero for pinned supports), we
solve for the unknown free-DOF displacements by reducing
\cref{eq:partition_system}:
\begin{equation}
  \boxed{\;
  \mathbf{K}_{\mathcal{F}\mathcal{F}}\,\mathbf{u}_{\mathcal{F}}
  \;=\;\mathbf{F}_{\mathcal{F}}-\mathbf{K}_{\mathcal{F}\mathcal{S}}\mathbf{u}_{\mathcal{S}}
  \;}
  \label{eq:reduced_system}
\end{equation}
which is a well-posed symmetric positive-definite linear system
provided the supports remove all rigid-body modes. For the
standard zero-prescribed-displacement case (a pinned or roller
support) \cref{eq:reduced_system} collapses to
$\mathbf{K}_{\mathcal{F}\mathcal{F}}\mathbf{u}_{\mathcal{F}}=\mathbf{F}_{\mathcal{F}}$.

The minimum number of DOFs to pin is
\begin{equation}
  |\mathcal{S}|_{\min}\;=\;\tfrac{d(d+1)}{2}
  \;=\;\begin{cases}3 & d=2\\6 & d=3\end{cases}.
\end{equation}
A 2D truss therefore needs at least one pin (two DOFs) plus one
roller (one DOF); a 3D truss needs a full pin and enough
additional roller DOFs to kill the three rotations. In our code
(\texttt{src/fem/solver.py:88}) failure to pin this minimum is
caught explicitly, and an under-constrained structure yields the
\texttt{LinAlgError: Singular matrix} that
\cref{eq:reduced_system} would otherwise produce only after a
solver attempt.

\paragraph{Reactions.} Once $\mathbf{u}_{\mathcal{F}}$ is known,
the reactions at supports follow from the second row of
\cref{eq:partition_system}:
\begin{equation}
  \mathbf{R}_{\mathcal{S}}
  \;=\;\mathbf{K}_{\mathcal{S}\mathcal{F}}\mathbf{u}_{\mathcal{F}}
  +\mathbf{K}_{\mathcal{S}\mathcal{S}}\mathbf{u}_{\mathcal{S}}
  -\mathbf{F}_{\mathcal{S}},
  \label{eq:reactions}
\end{equation}
which the solver returns alongside the displacements for
post-processing and as a global-equilibrium sanity check.

\section{Post-Processing: Stress, Axial Force, and Equilibrium}
\label{app:fem:postprocess}

After the reduced system is solved, the full displacement vector
$\mathbf{u}\in\mathbb{R}^{Nd}$ is known; stresses, strains and
member forces are recovered element by element:
\begin{enumerate}
  \item \textbf{Extract element displacements.} For element $e$
    pick out the $2d$ components
    $\mathbf{u}^{(e)}=\mathbf{u}[\mathcal{I}^{(e)}]$.
  \item \textbf{Axial strain} from
    \cref{eq:strain_displacement}:
    $\varepsilon^{(e)}=\mathbf{B}^{(e)}\mathbf{u}^{(e)}
    =\hat{\mathbf{n}}^{(e)}\!\cdot\!(\mathbf{u}_j-\mathbf{u}_i)/L^{(e)}$.
  \item \textbf{Axial stress} from \cref{eq:hooke1d}:
    $\sigma^{(e)}=E\,\varepsilon^{(e)}$.
  \item \textbf{Axial force} from \cref{eq:axial_force}:
    $N^{(e)}=A^{(e)}\sigma^{(e)}$.
\end{enumerate}
All four quantities are scalars; their sign convention is
``tension positive''.

\paragraph{Global equilibrium check.} A good implementation
verifies at post-processing time that
\begin{equation}
  \sum_{p\in\mathcal{S}}\mathbf{R}_p + \sum_{p=1}^{N}\mathbf{f}_p
  \;=\;\mathbf{0}
  \label{eq:global_equilibrium}
\end{equation}
to machine precision. Violations flag integration mistakes in
the load vector, forgotten supports, or incorrectly scattered
element contributions. Our unit test
\texttt{tests/test\_fem\_3bar.py} includes this check.

\paragraph{Linear-elastic assumptions and their scope.} The
derivation above is exact under H1--H3. In the thesis
benchmarks, H1 and H2 are introduced axiomatically by the
problem statements (every published truss optimum of
\cref{chap:results} uses the same assumptions). H3 is the only
one that can be challenged post-hoc: non-linearity shows up in
two forms:
\begin{itemize}
  \item \textbf{Material non-linearity}: yielding in steel beyond
    the elastic limit $f_y$. This is \emph{not} captured by
    \cref{eq:hooke1d}. For the IS~800 compliance study of
    \cref{sec:results:is800}, the tension-yield stress limit
    is enforced as an \emph{inequality constraint}
    $|\sigma^{(e)}|\le f_y$, not as a change of constitutive
    law, so linear elasticity remains the valid operating
    assumption throughout the optimiser's search space.
  \item \textbf{Geometric non-linearity}: buckling of slender
    compression members. \Cref{eq:strong_equilibrium} is the
    first-order form; it does not predict bifurcation. We
    therefore do not use the FEM to decide buckling: the IS~800
    buckling curve gives a closed-form compression capacity
    $P_d$ and the optimiser enforces $|N^{(e)}|\le P_d$ directly
    on element compression forces
    (\cref{sec:results:is800}).
\end{itemize}
Both exclusions are standard in the benchmark literature
~\cite{subramanian2008,is800-2007} and are faithfully reflected
in the constraint module \texttt{src/constraints/is800\_checks.py}.

\section{End-to-End Worked Example}
\label{app:fem:worked}

To ground the general derivation, we work through the single
horizontal bar in \texttt{src/fem/solver.py:147}. A bar of length
$L=\SI{1}{m}$ connects joint~$1$ at $(0,0)$ and joint~$2$ at
$(1,0)$, with $E=\SI{200}{GPa}$ and
$A=\SI{1e-4}{m^2}$. Joint~$1$ is fully pinned ($\mathcal{S}\supset\{1,2\}$)
and joint~$2$ carries an in-plane horizontal load $F_{2x}=\SI{1000}{N}$.
The $y$-DOF of joint~$2$ (DOF index $4$) must also be pinned to
remove the transverse zero mode of a single bar in 2D.

\emph{Step 1 --- Direction cosines and $\mathbf{B}^{(e)}$:}
$\hat{\mathbf{n}}=(1,0)^\top$, so
$\mathbf{B}^{(e)}=L^{-1}(-1,0,1,0)=(-1,0,1,0)$.

\emph{Step 2 --- Element stiffness:}
\[
  \mathbf{k}^{(e)}
  =\frac{EA}{L}\mathbf{B}^{(e)\top}\mathbf{B}^{(e)}
  =\SI{2e7}{N/m}\begin{pmatrix}
    +1 & 0 & -1 & 0\\ 0 & 0 & 0 & 0\\
    -1 & 0 & +1 & 0\\ 0 & 0 & 0 & 0
  \end{pmatrix}.
\]
With only one element $\mathbf{K}=\mathbf{k}^{(e)}$.

\emph{Step 3 --- Partition:} the free set is
$\mathcal{F}=\{3\}$ (the horizontal DOF of joint~$2$); the
prescribed set is $\mathcal{S}=\{1,2,4\}$ with
$\mathbf{u}_{\mathcal{S}}=\mathbf{0}$.
$\mathbf{K}_{\mathcal{F}\mathcal{F}}=\SI{2e7}{N/m}$,
$\mathbf{F}_{\mathcal{F}}=\SI{1000}{N}$.

\emph{Step 4 --- Solve reduced system}
\cref{eq:reduced_system}:
$u_{2x}=\SI{1000}{N}/\SI{2e7}{N/m}=\SI{5e-5}{m}=\SI{50}{\micro\metre}$.

\emph{Step 5 --- Axial strain, stress, force:}
$\varepsilon=\SI{5e-5}{\per\metre}\cdot L=\num{5e-5}$,
$\sigma=E\varepsilon=\SI{10}{MPa}$,
$N=A\sigma=\SI{1000}{N}$ (tension).

\emph{Step 6 --- Reaction at joint 1:} from
\cref{eq:reactions}, $R_{1x}=-\SI{1000}{N}$, satisfying
global equilibrium to machine precision.

Every step above appears identically in
\texttt{src/fem/solver.py}'s
\texttt{if \_\_name\_\_==\"\_\_main\_\_\"} block, and its
assertion $u_{2x}=\SI{5e-5}{m}$ is the simplest validation gate
of the FEM stack.

\section{Notation Summary}
\label{app:fem:notation}

\begin{tabular}{@{}l l l@{}}
  \toprule
  Symbol & Meaning & Units \\
  \midrule
  $d$ & spatial dimension (2 or 3) & --- \\
  $N$ & number of joints & --- \\
  $M$ & number of bars & --- \\
  $E$ & Young's modulus of steel & Pa \\
  $A^{(e)}$ & cross-sectional area of bar $e$ & $\mathrm{m}^2$ \\
  $L^{(e)}$ & undeformed length of bar $e$ & m \\
  $\hat{\mathbf{n}}^{(e)}$ & unit direction vector of bar $e$ & --- \\
  $\mathbf{B}^{(e)}$ & strain-displacement matrix of bar $e$ ($1\!\times\!2d$) & $\mathrm{m}^{-1}$ \\
  $\mathbf{k}^{(e)}$ & element stiffness matrix ($2d\!\times\!2d$) & N/m \\
  $\mathbf{K}$ & global stiffness matrix ($Nd\!\times\!Nd$) & N/m \\
  $\mathbf{u}$ & nodal-displacement vector & m \\
  $\mathbf{F}$ & nodal-force vector & N \\
  $\varepsilon^{(e)},\sigma^{(e)},N^{(e)}$ & element axial strain, stress, force & ---, Pa, N \\
  $\mathcal{F},\mathcal{S}$ & free / prescribed DOF sets & --- \\
  \bottomrule
\end{tabular}
